title: User-Configurable Settings
alias: settings, configuration
todo:
------------------------------
\usepackage{listings}
\lstset{language=C++,basicstyle=\small\ttfamily,backgroundcolor=\color{black!10!white},columns=fullflexible,tabsize=4}
------------------------------

Every user-configurable value in GDSiMS lives in a single settings object, constructed from [[lsj|reflection of the JSL configuration module]]. This gives each setting a command-line flag, a configuration-file entry, an entry in the help screen -- amongst others -- with no more code than declaring the variable itself.

\section{Structure}

	The settings are grouped into sections, each a class in _src/settings/_, which are collected together in _MainSettings_:

	* _Runtime_: how the program runs (logging verbosity, number of threads), as opposed to what it simulates.
	* _Simulation_: the parameters of the simulation itself.

	A section can itself contain further sections, and the parsing, help and configuration all descend through them automatically.

\section{Setting values}

	A setting can be given a value in three ways, in increasing order of priority:

	* Its default, written in the class declaration.
	* A configuration file, given with _-config file_, containing one _key value_ pair per line.
	* The command line.

	So, for example, a configuration file can hold a standard set of parameters, with one or two then changed on the command line for a particular run:

	\begin{lstlisting}[language=bash]
./gdsims -config standard.cfg -T 365
\end{lstlisting}

	The full list of settings, with their flags and defaults, is always available from _./gdsims -h_.

\section{Accessing the settings in code}

	There is a single global settings object, _Settings_. For all intents and purposes, this is a **read-only** object, which cannot be modified. However, since it is global, the value of the settings can be read by any file.

	\subsection{Changing the Settings}

		***Do Not Do This***.

		The settings are only changed when the user directs their value be changed (such as at startup or during an interactive session) -- never during the normal running of the code.

\section{Derived values}

	Some things are not settings in their own right, but follow from them: for example, the log level follows from the _-q_, _-s_ and _-v_ flags. These are computed in a single place, _deriveSettingsValues()_ (_src/settings/derive.cpp_), which is called every time the settings change. Anything which is calculated from the settings, rather than set directly, belongs there.

\section{Adding a setting}

	Add a documented member to the relevant section:

	\begin{lstlisting}
//! @brief The number of timesteps to run the simulation for
//! @alias T, TimeSteps
int TimeSteps = 100;
\end{lstlisting}

	Then rebuild: the [[lsj]] mirror runs as part of the build, and the new setting immediately has a flag (any of the assigned aliases, such as _-T_), a help entry and a configuration-file key.

	If the setting has a constraint (a value which must be positive, or two settings which are incompatible), add a check for it to the [[validation]].
